\section{Objetivos}

\subsection{Objetivo General}

Diseñar y desarrollar un sistema de gestión y trazabilidad de servicios de soporte técnico con asistencia mediante inteligencia artificial para los prestadores de servicios de la ciudad de Tarija durante las gestiones 2026--2027, a partir del análisis de los procesos actuales y del levantamiento de requerimientos, con la finalidad de centralizar el registro de las órdenes de trabajo y facilitar su seguimiento y la comunicación con los clientes.

\subsection{Objetivos Específicos}

\begin{enumerate}
    \item Analizar el proceso actual de recepción de equipos, registro de servicios y comunicación con los clientes en prestadores de soporte técnico de la ciudad de Tarija, identificando las principales necesidades de información.
    \item Determinar los requerimientos funcionales y no funcionales del sistema, considerando el registro de datos de recepción, la gestión de estados del servicio y las necesidades de consulta de los usuarios.
    \item Modelar la estructura del sistema, los flujos de trabajo de las órdenes de servicio y la arquitectura del software orientada al desarrollo de la solución.
    \item Diseñar la lógica de interacción de un asistente conversacional basado en inteligencia artificial para apoyar la atención de consultas frecuentes y la entrega de información sobre el estado de los servicios.
    \item Desarrollar el prototipo funcional (MVP) de la solución propuesta y evaluar su factibilidad operativa mediante criterios y métricas de calidad de software según la norma ISO/IEC 25010:2023 para validar el sistema de trazabilidad de soporte técnico.
\end{enumerate}